\section{高道研究室・稲垣賢斗さんの発表について}

\subsection{Audio Captioningモデルの発達的カリキュラム学習}
本発表では，幼児の言語発達過程に着想を得た「発達的カリキュラム学習」によって，Audio Captioningモデルの性能向上を目指す手法が提案された．語彙数・文法複雑性・文長などを段階的に変化させた学習データを用いる点が特徴である．

\subsubsection{質問の背景と意図}
訓練データは年齢段階を模した加工がされていたが，評価時に使われる正解キャプションがどの段階に対応するのかが気になった．特に，文長が比較的長く感じられたため，それが「大人の発話」に相当するものなのかを確認したかった．

\subsubsection{質問内容}
正解文はカリキュラム上どの段階に相当するのか？文長30という観点からも，それは「最終段階＝大人」レベルでよいか？

\subsubsection{回答内容}
正解文はカリキュラムにおける最終段階に相当し，大人の人間が書いたtestセットのキャプションである．文の長さは固定ではなく，音の内容に依存してまちまちとのことだった．

\subsubsection{考察}
モデルには簡単な文から徐々に学習させる一方で，評価には最終目標である大人文を使っている点が納得できた．評価軸として自然文に基づく一貫した指標があることで，訓練プロセスと性能の整合性が保たれていると感じた．